% =============================================================================
% APPENDIX: Technical details of Theorem 1 (Cosmological constant subsection)
% Included via % =============================================================================
% APPENDIX: Technical details of Theorem 1 (Cosmological constant subsection)
% Included via % =============================================================================
% APPENDIX: Technical details of Theorem 1 (Cosmological constant subsection)
% Included via % =============================================================================
% APPENDIX: Technical details of Theorem 1 (Cosmological constant subsection)
% Included via \input{App_CC_technical} in ECT_preprint.tex
% =============================================================================

\section{Technical details for the emergent unimodular decoupling
theorem}
\label{app:cc_technical}

This appendix collects the explicit derivations underlying the
cosmological constant analysis of Section~\ref{sec:cc_problem}. We
give (i) the rank-one determinant identity with the admissible-variation
argument; (ii) the Jacobian cancellation identity for the classical
baseline; (iii) the FRW reduction leading to the slow-branch equation
of state; and (iv) the dimensional-consistency audit of the infrared
scaling with the pseudo-Goldstone candidate.


\subsection{Rank-one determinant identity and fixed-determinant
property}
\label{app:rank1}

The matrix determinant lemma states, for a symmetric $4\times 4$
matrix $M$ and a 4-vector $n$,
\begin{equation}
\det(M + c\,n n^{\mathsf T})
\;=\; (1 + c\,n^{\mathsf T} M^{-1} n)\,\det M\,,
\label{eq:matrix_det_lemma}
\end{equation}
provided $M$ is invertible. Specialising to $M = \beta\,\delta$ and
$c = -\alpha$:
\begin{align}
\det(\beta\,\delta^{AB} - \alpha\,n^A n^B)
&\;=\; \beta^4 \cdot \Bigl(1 - \tfrac{\alpha}{\beta}\,n^A
\delta_{AB} n^B\Bigr)\,.
\label{eq:K_det_general_app}
\end{align}

In the ordered branch the orientation field satisfies the
constraint $n^A(x)\,\delta_{AB}\,n^B(x) = 1$, reflecting that $|n|$
is not an independent degree of freedom (amplitude variation is
described by the radial mode $\phi$). Substituting:
\begin{equation}
\det K^{AB} \;=\; \beta^4\left(1 - \tfrac{\alpha}{\beta}\right)
\;=\; \beta^3(\beta - \alpha) \;=\; -\beta^3(\alpha - \beta)\,,
\label{eq:detK_app}
\end{equation}
independent of $n^A(x)$.

From the densitized-metric identification $K^{AB} \propto
\sqrt{-g_{\mathrm{eff}}}\,g_{\mathrm{eff}}^{AB}$ and
$\det(\sqrt{-g}\,g^{\mu\nu}) = -g$ in four dimensions,
\begin{equation}
\det K^{AB} \;=\; -\,g_{\mathrm{eff}}\,,
\qquad
\sqrt{-g_{\mathrm{eff}}} \;=\;
\sqrt{\beta^3(\alpha-\beta)} \;=\; \frac{\beta^2}{c_*}\,,
\label{eq:sqrt_g_value_app}
\end{equation}
using $c_*^2 = \beta/(\alpha-\beta)$.

\paragraph{Vanishing variation on the constraint surface.} An
admissible variation $\delta n^A$ preserves the normalization
$n^A\delta_{AB}n^B=1$, hence $\delta_{AB} n^A \delta n^B = 0$. Differentiating
\eqref{eq:K_det_general_app},
\begin{equation}
\frac{\partial(\det K)}{\partial n^A}
\;=\; -2\alpha\beta^3\, \delta_{AB} n^B\,,
\end{equation}
hence
\begin{equation}
\delta(\det K)
\;=\; -2\alpha\beta^3\,(\delta_{AB} n^B \delta n^A) \;=\; 0\,,
\end{equation}
so $\delta\sqrt{-g_{\mathrm{eff}}} = 0$ on admissible variations, as
stated in Theorem~1(iii) of the main text.


\subsection{Jacobian cancellation identity for the classical
baseline}
\label{app:jacobian_cc}

The fundamental ECT action is $S_E[q] = \int_M d^4 x_E\,
\mathcal{L}_E(q,\partial q)$ with $q^a$ collectively denoting
$\Phi$, $n^A$, and auxiliary fields. Crucially, the emergent metric
$g_{\mathrm{eff}}^{AB}[q]$ is a functional of the condensate fields
rather than an independent dynamical variable; there is no variation
step ``by $g_{\mu\nu}$''.

Under the shift $\mathcal{L}_E \to \mathcal{L}_E + C$ for a
field-independent constant $C\in\mathbb{R}$:
(a) $\partial C/\partial q^a = \partial C/\partial(\partial_A q^a)
= 0$, so every Euler--Lagrange equation is unchanged; (b) the
classical solutions $q^a(x)$ are unchanged; (c) the emergent metric
$g_{\mathrm{eff}}^{AB}[q(x)]$ is unchanged on each classical
solution. Specialising to $C = V(\phi_0)$ establishes Theorem~1(iv).

The result is non-trivial in the ECT context precisely because
$g_{\mathrm{eff}}$ is emergent. In standard general relativity with
an independent $g_{\mu\nu}$, the same shift generates $C\sqrt{-g}$,
whose variation produces $\sim C g_{\mu\nu}\sqrt{-g}$ -- the standard
cosmological-constant source. That mechanism is unavailable here.


\subsection{FRW reduction and equation of state}
\label{app:frw_cc}

On an FRW background, spatial isotropy selects a single
long-wavelength ordered-branch degree of freedom $q(t)$. The
minimal homogeneous effective Lagrangian is
\begin{equation}
L_{\mathrm{IR}}[q; a] \;=\; a^3 \left[
\tfrac{1}{2} Z_q\, \dot q^2 - V_{\mathrm{IR}}(q) \right],
\label{eq:frw_lagrangian_app}
\end{equation}
with infrared stiffness $Z_q$ and residual infrared potential
$V_{\mathrm{IR}}(q)$. Crucially, $V_{\mathrm{IR}}(q)$ is the
dynamical dressing of the potential around the ordered-branch
minimum; it is not $V(\phi_0)$, which is decoupled by Theorem~1(iv).
This distinction resolves the apparent tension between the
decoupling theorem and the dark-energy picture used in
Sec.~\ref{sec:mp_late_cosmology}.

From \eqref{eq:frw_lagrangian_app},
\begin{align}
\rho_q &\;=\; \tfrac{1}{2} Z_q \dot q^2 + V_{\mathrm{IR}}(q)\,,
\\
p_q &\;=\; \tfrac{1}{2} Z_q \dot q^2 - V_{\mathrm{IR}}(q)\,,
\\
w_q &\;=\; \frac{r - 1}{r + 1}\,,
\qquad r \equiv
\frac{\tfrac12 Z_q \dot q^2}{V_{\mathrm{IR}}(q)}\,.
\end{align}
For $r \ll 1$ (slow-branch regime), $w_q \simeq -1 + 2r$,
reproducing the main-text formula $w_0 = -1 +
2\rho_{\mathrm{kin}}/(3\rho_{\mathrm{cond}})$ upon identification
$\rho_{\mathrm{kin}}=\tfrac12 Z_q\dot q^2$ and
$\rho_{\mathrm{cond}}=V_{\mathrm{IR}}(q)$.

\subsection{Dimensional audit of the infrared scaling}
\label{app:ir_audit_cc}

The observed $\rho_\Lambda$ has dimensions mass$^4$. After blocking
the direct UV channel (Theorem~1), the scale $M_{\mathrm{Pl}}^4$ is
not available as a direct source. Available infrared scales are the
emergent gravitational stiffness $M_{\mathrm{Pl}}^2$ (from the
Sakharov/Adler induced-gravity mechanism~\cite{Sakharov1967,Adler1982})
and the Hubble scale $H_0$ (only dimensionful IR curvature/time
scale).

Dimensionally consistent combinations of mass$^4$ from
$\{M_{\mathrm{Pl}}^2, H_0\}$ without directly invoking
$M_{\mathrm{Pl}}^4$ are $M_{\mathrm{Pl}}^2 H_0^2$, $M_{\mathrm{Pl}}
H_0^3$, and $H_0^4$. With $M_{\mathrm{Pl}}/H_0 \sim 10^{60}$, these
evaluate to $10^{-47}$, $10^{-107}$, $10^{-167}\,\mathrm{GeV}^4$
respectively; only the first matches observation at the
order-of-magnitude level.

\paragraph{Pseudo-Goldstone realization.} For a pseudo-Goldstone
$\theta$ of mass $m_G$ and decay constant $f_G$, the natural infrared
energy density is $\rho_G \sim m_G^2 f_G^2$, obtained from
$V_G(\theta) \sim m_G^2\theta^2$ at $\theta \sim f_G$. Using the ECT
values $m_G \sim H_0$ and $f_G \sim \phi_0 \sim M_{\mathrm{Pl}}$
(Sec.~\ref{subsec:goldstone_cc}),
\begin{equation}
\rho_G \;\sim\; m_G^2 f_G^2 \;\sim\; H_0^2 M_{\mathrm{Pl}}^2\,,
\end{equation}
reproducing the IR scaling with $c_\Lambda = O(1)$ via independent
inputs: $m_G$ comes from the pseudo-Goldstone (dark-matter-candidate)
sector, $f_G$ from the gravitational SSB scale. Neither is tuned
here.

\paragraph{Scope.} The estimate $\rho_G \sim H_0^2 M_{\mathrm{Pl}}^2$
is a consistency check, not a derivation. Upgrading it to a
prediction requires a first-principles derivation of $m_G$ from ECT
parameters (Open Problem~O4) and an analysis of the cosmological
evolution of the Goldstone condensate. The holographic bound of
Cohen--Kaplan--Nelson~\cite{CohenKaplanNelson1999} states
$\rho \lesssim M_{\mathrm{Pl}}^2/L_{\mathrm{IR}}^2$; the ECT estimate
saturates this bound with $L_{\mathrm{IR}}\sim H_0^{-1}$, which is an
independent phenomenological consistency check.

% =============================================================================
% END OF APPENDIX App_CC_technical
% ============================================================================= in ECT_preprint.tex
% =============================================================================

\section{Technical details for the emergent unimodular decoupling
theorem}
\label{app:cc_technical}

This appendix collects the explicit derivations underlying the
cosmological constant analysis of Section~\ref{sec:cc_problem}. We
give (i) the rank-one determinant identity with the admissible-variation
argument; (ii) the Jacobian cancellation identity for the classical
baseline; (iii) the FRW reduction leading to the slow-branch equation
of state; and (iv) the dimensional-consistency audit of the infrared
scaling with the pseudo-Goldstone candidate.


\subsection{Rank-one determinant identity and fixed-determinant
property}
\label{app:rank1}

The matrix determinant lemma states, for a symmetric $4\times 4$
matrix $M$ and a 4-vector $n$,
\begin{equation}
\det(M + c\,n n^{\mathsf T})
\;=\; (1 + c\,n^{\mathsf T} M^{-1} n)\,\det M\,,
\label{eq:matrix_det_lemma}
\end{equation}
provided $M$ is invertible. Specialising to $M = \beta\,\delta$ and
$c = -\alpha$:
\begin{align}
\det(\beta\,\delta^{AB} - \alpha\,n^A n^B)
&\;=\; \beta^4 \cdot \Bigl(1 - \tfrac{\alpha}{\beta}\,n^A
\delta_{AB} n^B\Bigr)\,.
\label{eq:K_det_general_app}
\end{align}

In the ordered branch the orientation field satisfies the
constraint $n^A(x)\,\delta_{AB}\,n^B(x) = 1$, reflecting that $|n|$
is not an independent degree of freedom (amplitude variation is
described by the radial mode $\phi$). Substituting:
\begin{equation}
\det K^{AB} \;=\; \beta^4\left(1 - \tfrac{\alpha}{\beta}\right)
\;=\; \beta^3(\beta - \alpha) \;=\; -\beta^3(\alpha - \beta)\,,
\label{eq:detK_app}
\end{equation}
independent of $n^A(x)$.

From the densitized-metric identification $K^{AB} \propto
\sqrt{-g_{\mathrm{eff}}}\,g_{\mathrm{eff}}^{AB}$ and
$\det(\sqrt{-g}\,g^{\mu\nu}) = -g$ in four dimensions,
\begin{equation}
\det K^{AB} \;=\; -\,g_{\mathrm{eff}}\,,
\qquad
\sqrt{-g_{\mathrm{eff}}} \;=\;
\sqrt{\beta^3(\alpha-\beta)} \;=\; \frac{\beta^2}{c_*}\,,
\label{eq:sqrt_g_value_app}
\end{equation}
using $c_*^2 = \beta/(\alpha-\beta)$.

\paragraph{Vanishing variation on the constraint surface.} An
admissible variation $\delta n^A$ preserves the normalization
$n^A\delta_{AB}n^B=1$, hence $\delta_{AB} n^A \delta n^B = 0$. Differentiating
\eqref{eq:K_det_general_app},
\begin{equation}
\frac{\partial(\det K)}{\partial n^A}
\;=\; -2\alpha\beta^3\, \delta_{AB} n^B\,,
\end{equation}
hence
\begin{equation}
\delta(\det K)
\;=\; -2\alpha\beta^3\,(\delta_{AB} n^B \delta n^A) \;=\; 0\,,
\end{equation}
so $\delta\sqrt{-g_{\mathrm{eff}}} = 0$ on admissible variations, as
stated in Theorem~1(iii) of the main text.


\subsection{Jacobian cancellation identity for the classical
baseline}
\label{app:jacobian_cc}

The fundamental ECT action is $S_E[q] = \int_M d^4 x_E\,
\mathcal{L}_E(q,\partial q)$ with $q^a$ collectively denoting
$\Phi$, $n^A$, and auxiliary fields. Crucially, the emergent metric
$g_{\mathrm{eff}}^{AB}[q]$ is a functional of the condensate fields
rather than an independent dynamical variable; there is no variation
step ``by $g_{\mu\nu}$''.

Under the shift $\mathcal{L}_E \to \mathcal{L}_E + C$ for a
field-independent constant $C\in\mathbb{R}$:
(a) $\partial C/\partial q^a = \partial C/\partial(\partial_A q^a)
= 0$, so every Euler--Lagrange equation is unchanged; (b) the
classical solutions $q^a(x)$ are unchanged; (c) the emergent metric
$g_{\mathrm{eff}}^{AB}[q(x)]$ is unchanged on each classical
solution. Specialising to $C = V(\phi_0)$ establishes Theorem~1(iv).

The result is non-trivial in the ECT context precisely because
$g_{\mathrm{eff}}$ is emergent. In standard general relativity with
an independent $g_{\mu\nu}$, the same shift generates $C\sqrt{-g}$,
whose variation produces $\sim C g_{\mu\nu}\sqrt{-g}$ -- the standard
cosmological-constant source. That mechanism is unavailable here.


\subsection{FRW reduction and equation of state}
\label{app:frw_cc}

On an FRW background, spatial isotropy selects a single
long-wavelength ordered-branch degree of freedom $q(t)$. The
minimal homogeneous effective Lagrangian is
\begin{equation}
L_{\mathrm{IR}}[q; a] \;=\; a^3 \left[
\tfrac{1}{2} Z_q\, \dot q^2 - V_{\mathrm{IR}}(q) \right],
\label{eq:frw_lagrangian_app}
\end{equation}
with infrared stiffness $Z_q$ and residual infrared potential
$V_{\mathrm{IR}}(q)$. Crucially, $V_{\mathrm{IR}}(q)$ is the
dynamical dressing of the potential around the ordered-branch
minimum; it is not $V(\phi_0)$, which is decoupled by Theorem~1(iv).
This distinction resolves the apparent tension between the
decoupling theorem and the dark-energy picture used in
Sec.~\ref{sec:mp_late_cosmology}.

From \eqref{eq:frw_lagrangian_app},
\begin{align}
\rho_q &\;=\; \tfrac{1}{2} Z_q \dot q^2 + V_{\mathrm{IR}}(q)\,,
\\
p_q &\;=\; \tfrac{1}{2} Z_q \dot q^2 - V_{\mathrm{IR}}(q)\,,
\\
w_q &\;=\; \frac{r - 1}{r + 1}\,,
\qquad r \equiv
\frac{\tfrac12 Z_q \dot q^2}{V_{\mathrm{IR}}(q)}\,.
\end{align}
For $r \ll 1$ (slow-branch regime), $w_q \simeq -1 + 2r$,
reproducing the main-text formula $w_0 = -1 +
2\rho_{\mathrm{kin}}/(3\rho_{\mathrm{cond}})$ upon identification
$\rho_{\mathrm{kin}}=\tfrac12 Z_q\dot q^2$ and
$\rho_{\mathrm{cond}}=V_{\mathrm{IR}}(q)$.

\subsection{Dimensional audit of the infrared scaling}
\label{app:ir_audit_cc}

The observed $\rho_\Lambda$ has dimensions mass$^4$. After blocking
the direct UV channel (Theorem~1), the scale $M_{\mathrm{Pl}}^4$ is
not available as a direct source. Available infrared scales are the
emergent gravitational stiffness $M_{\mathrm{Pl}}^2$ (from the
Sakharov/Adler induced-gravity mechanism~\cite{Sakharov1967,Adler1982})
and the Hubble scale $H_0$ (only dimensionful IR curvature/time
scale).

Dimensionally consistent combinations of mass$^4$ from
$\{M_{\mathrm{Pl}}^2, H_0\}$ without directly invoking
$M_{\mathrm{Pl}}^4$ are $M_{\mathrm{Pl}}^2 H_0^2$, $M_{\mathrm{Pl}}
H_0^3$, and $H_0^4$. With $M_{\mathrm{Pl}}/H_0 \sim 10^{60}$, these
evaluate to $10^{-47}$, $10^{-107}$, $10^{-167}\,\mathrm{GeV}^4$
respectively; only the first matches observation at the
order-of-magnitude level.

\paragraph{Pseudo-Goldstone realization.} For a pseudo-Goldstone
$\theta$ of mass $m_G$ and decay constant $f_G$, the natural infrared
energy density is $\rho_G \sim m_G^2 f_G^2$, obtained from
$V_G(\theta) \sim m_G^2\theta^2$ at $\theta \sim f_G$. Using the ECT
values $m_G \sim H_0$ and $f_G \sim \phi_0 \sim M_{\mathrm{Pl}}$
(Sec.~\ref{subsec:goldstone_cc}),
\begin{equation}
\rho_G \;\sim\; m_G^2 f_G^2 \;\sim\; H_0^2 M_{\mathrm{Pl}}^2\,,
\end{equation}
reproducing the IR scaling with $c_\Lambda = O(1)$ via independent
inputs: $m_G$ comes from the pseudo-Goldstone (dark-matter-candidate)
sector, $f_G$ from the gravitational SSB scale. Neither is tuned
here.

\paragraph{Scope.} The estimate $\rho_G \sim H_0^2 M_{\mathrm{Pl}}^2$
is a consistency check, not a derivation. Upgrading it to a
prediction requires a first-principles derivation of $m_G$ from ECT
parameters (Open Problem~O4) and an analysis of the cosmological
evolution of the Goldstone condensate. The holographic bound of
Cohen--Kaplan--Nelson~\cite{CohenKaplanNelson1999} states
$\rho \lesssim M_{\mathrm{Pl}}^2/L_{\mathrm{IR}}^2$; the ECT estimate
saturates this bound with $L_{\mathrm{IR}}\sim H_0^{-1}$, which is an
independent phenomenological consistency check.

% =============================================================================
% END OF APPENDIX App_CC_technical
% ============================================================================= in ECT_preprint.tex
% =============================================================================

\section{Technical details for the emergent unimodular decoupling
theorem}
\label{app:cc_technical}

This appendix collects the explicit derivations underlying the
cosmological constant analysis of Section~\ref{sec:cc_problem}. We
give (i) the rank-one determinant identity with the admissible-variation
argument; (ii) the Jacobian cancellation identity for the classical
baseline; (iii) the FRW reduction leading to the slow-branch equation
of state; and (iv) the dimensional-consistency audit of the infrared
scaling with the pseudo-Goldstone candidate.


\subsection{Rank-one determinant identity and fixed-determinant
property}
\label{app:rank1}

The matrix determinant lemma states, for a symmetric $4\times 4$
matrix $M$ and a 4-vector $n$,
\begin{equation}
\det(M + c\,n n^{\mathsf T})
\;=\; (1 + c\,n^{\mathsf T} M^{-1} n)\,\det M\,,
\label{eq:matrix_det_lemma}
\end{equation}
provided $M$ is invertible. Specialising to $M = \beta\,\delta$ and
$c = -\alpha$:
\begin{align}
\det(\beta\,\delta^{AB} - \alpha\,n^A n^B)
&\;=\; \beta^4 \cdot \Bigl(1 - \tfrac{\alpha}{\beta}\,n^A
\delta_{AB} n^B\Bigr)\,.
\label{eq:K_det_general_app}
\end{align}

In the ordered branch the orientation field satisfies the
constraint $n^A(x)\,\delta_{AB}\,n^B(x) = 1$, reflecting that $|n|$
is not an independent degree of freedom (amplitude variation is
described by the radial mode $\phi$). Substituting:
\begin{equation}
\det K^{AB} \;=\; \beta^4\left(1 - \tfrac{\alpha}{\beta}\right)
\;=\; \beta^3(\beta - \alpha) \;=\; -\beta^3(\alpha - \beta)\,,
\label{eq:detK_app}
\end{equation}
independent of $n^A(x)$.

From the densitized-metric identification $K^{AB} \propto
\sqrt{-g_{\mathrm{eff}}}\,g_{\mathrm{eff}}^{AB}$ and
$\det(\sqrt{-g}\,g^{\mu\nu}) = -g$ in four dimensions,
\begin{equation}
\det K^{AB} \;=\; -\,g_{\mathrm{eff}}\,,
\qquad
\sqrt{-g_{\mathrm{eff}}} \;=\;
\sqrt{\beta^3(\alpha-\beta)} \;=\; \frac{\beta^2}{c_*}\,,
\label{eq:sqrt_g_value_app}
\end{equation}
using $c_*^2 = \beta/(\alpha-\beta)$.

\paragraph{Vanishing variation on the constraint surface.} An
admissible variation $\delta n^A$ preserves the normalization
$n^A\delta_{AB}n^B=1$, hence $\delta_{AB} n^A \delta n^B = 0$. Differentiating
\eqref{eq:K_det_general_app},
\begin{equation}
\frac{\partial(\det K)}{\partial n^A}
\;=\; -2\alpha\beta^3\, \delta_{AB} n^B\,,
\end{equation}
hence
\begin{equation}
\delta(\det K)
\;=\; -2\alpha\beta^3\,(\delta_{AB} n^B \delta n^A) \;=\; 0\,,
\end{equation}
so $\delta\sqrt{-g_{\mathrm{eff}}} = 0$ on admissible variations, as
stated in Theorem~1(iii) of the main text.


\subsection{Jacobian cancellation identity for the classical
baseline}
\label{app:jacobian_cc}

The fundamental ECT action is $S_E[q] = \int_M d^4 x_E\,
\mathcal{L}_E(q,\partial q)$ with $q^a$ collectively denoting
$\Phi$, $n^A$, and auxiliary fields. Crucially, the emergent metric
$g_{\mathrm{eff}}^{AB}[q]$ is a functional of the condensate fields
rather than an independent dynamical variable; there is no variation
step ``by $g_{\mu\nu}$''.

Under the shift $\mathcal{L}_E \to \mathcal{L}_E + C$ for a
field-independent constant $C\in\mathbb{R}$:
(a) $\partial C/\partial q^a = \partial C/\partial(\partial_A q^a)
= 0$, so every Euler--Lagrange equation is unchanged; (b) the
classical solutions $q^a(x)$ are unchanged; (c) the emergent metric
$g_{\mathrm{eff}}^{AB}[q(x)]$ is unchanged on each classical
solution. Specialising to $C = V(\phi_0)$ establishes Theorem~1(iv).

The result is non-trivial in the ECT context precisely because
$g_{\mathrm{eff}}$ is emergent. In standard general relativity with
an independent $g_{\mu\nu}$, the same shift generates $C\sqrt{-g}$,
whose variation produces $\sim C g_{\mu\nu}\sqrt{-g}$ -- the standard
cosmological-constant source. That mechanism is unavailable here.


\subsection{FRW reduction and equation of state}
\label{app:frw_cc}

On an FRW background, spatial isotropy selects a single
long-wavelength ordered-branch degree of freedom $q(t)$. The
minimal homogeneous effective Lagrangian is
\begin{equation}
L_{\mathrm{IR}}[q; a] \;=\; a^3 \left[
\tfrac{1}{2} Z_q\, \dot q^2 - V_{\mathrm{IR}}(q) \right],
\label{eq:frw_lagrangian_app}
\end{equation}
with infrared stiffness $Z_q$ and residual infrared potential
$V_{\mathrm{IR}}(q)$. Crucially, $V_{\mathrm{IR}}(q)$ is the
dynamical dressing of the potential around the ordered-branch
minimum; it is not $V(\phi_0)$, which is decoupled by Theorem~1(iv).
This distinction resolves the apparent tension between the
decoupling theorem and the dark-energy picture used in
Sec.~\ref{sec:mp_late_cosmology}.

From \eqref{eq:frw_lagrangian_app},
\begin{align}
\rho_q &\;=\; \tfrac{1}{2} Z_q \dot q^2 + V_{\mathrm{IR}}(q)\,,
\\
p_q &\;=\; \tfrac{1}{2} Z_q \dot q^2 - V_{\mathrm{IR}}(q)\,,
\\
w_q &\;=\; \frac{r - 1}{r + 1}\,,
\qquad r \equiv
\frac{\tfrac12 Z_q \dot q^2}{V_{\mathrm{IR}}(q)}\,.
\end{align}
For $r \ll 1$ (slow-branch regime), $w_q \simeq -1 + 2r$,
reproducing the main-text formula $w_0 = -1 +
2\rho_{\mathrm{kin}}/(3\rho_{\mathrm{cond}})$ upon identification
$\rho_{\mathrm{kin}}=\tfrac12 Z_q\dot q^2$ and
$\rho_{\mathrm{cond}}=V_{\mathrm{IR}}(q)$.

\subsection{Dimensional audit of the infrared scaling}
\label{app:ir_audit_cc}

The observed $\rho_\Lambda$ has dimensions mass$^4$. After blocking
the direct UV channel (Theorem~1), the scale $M_{\mathrm{Pl}}^4$ is
not available as a direct source. Available infrared scales are the
emergent gravitational stiffness $M_{\mathrm{Pl}}^2$ (from the
Sakharov/Adler induced-gravity mechanism~\cite{Sakharov1967,Adler1982})
and the Hubble scale $H_0$ (only dimensionful IR curvature/time
scale).

Dimensionally consistent combinations of mass$^4$ from
$\{M_{\mathrm{Pl}}^2, H_0\}$ without directly invoking
$M_{\mathrm{Pl}}^4$ are $M_{\mathrm{Pl}}^2 H_0^2$, $M_{\mathrm{Pl}}
H_0^3$, and $H_0^4$. With $M_{\mathrm{Pl}}/H_0 \sim 10^{60}$, these
evaluate to $10^{-47}$, $10^{-107}$, $10^{-167}\,\mathrm{GeV}^4$
respectively; only the first matches observation at the
order-of-magnitude level.

\paragraph{Pseudo-Goldstone realization.} For a pseudo-Goldstone
$\theta$ of mass $m_G$ and decay constant $f_G$, the natural infrared
energy density is $\rho_G \sim m_G^2 f_G^2$, obtained from
$V_G(\theta) \sim m_G^2\theta^2$ at $\theta \sim f_G$. Using the ECT
values $m_G \sim H_0$ and $f_G \sim \phi_0 \sim M_{\mathrm{Pl}}$
(Sec.~\ref{subsec:goldstone_cc}),
\begin{equation}
\rho_G \;\sim\; m_G^2 f_G^2 \;\sim\; H_0^2 M_{\mathrm{Pl}}^2\,,
\end{equation}
reproducing the IR scaling with $c_\Lambda = O(1)$ via independent
inputs: $m_G$ comes from the pseudo-Goldstone (dark-matter-candidate)
sector, $f_G$ from the gravitational SSB scale. Neither is tuned
here.

\paragraph{Scope.} The estimate $\rho_G \sim H_0^2 M_{\mathrm{Pl}}^2$
is a consistency check, not a derivation. Upgrading it to a
prediction requires a first-principles derivation of $m_G$ from ECT
parameters (Open Problem~O4) and an analysis of the cosmological
evolution of the Goldstone condensate. The holographic bound of
Cohen--Kaplan--Nelson~\cite{CohenKaplanNelson1999} states
$\rho \lesssim M_{\mathrm{Pl}}^2/L_{\mathrm{IR}}^2$; the ECT estimate
saturates this bound with $L_{\mathrm{IR}}\sim H_0^{-1}$, which is an
independent phenomenological consistency check.

% =============================================================================
% END OF APPENDIX App_CC_technical
% ============================================================================= in ECT_preprint.tex
% =============================================================================

\section{Technical details for the emergent unimodular decoupling
theorem}
\label{app:cc_technical}

This appendix collects the explicit derivations underlying the
cosmological constant analysis of Section~\ref{sec:cc_problem}. We
give (i) the rank-one determinant identity with the admissible-variation
argument; (ii) the Jacobian cancellation identity for the classical
baseline; (iii) the FRW reduction leading to the slow-branch equation
of state; and (iv) the dimensional-consistency audit of the infrared
scaling with the pseudo-Goldstone candidate.


\subsection{Rank-one determinant identity and fixed-determinant
property}
\label{app:rank1}

The matrix determinant lemma states, for a symmetric $4\times 4$
matrix $M$ and a 4-vector $n$,
\begin{equation}
\det(M + c\,n n^{\mathsf T})
\;=\; (1 + c\,n^{\mathsf T} M^{-1} n)\,\det M\,,
\label{eq:matrix_det_lemma}
\end{equation}
provided $M$ is invertible. Specialising to $M = \beta\,\delta$ and
$c = -\alpha$:
\begin{align}
\det(\beta\,\delta^{AB} - \alpha\,n^A n^B)
&\;=\; \beta^4 \cdot \Bigl(1 - \tfrac{\alpha}{\beta}\,n^A
\delta_{AB} n^B\Bigr)\,.
\label{eq:K_det_general_app}
\end{align}

In the ordered branch the orientation field satisfies the
constraint $n^A(x)\,\delta_{AB}\,n^B(x) = 1$, reflecting that $|n|$
is not an independent degree of freedom (amplitude variation is
described by the radial mode $\phi$). Substituting:
\begin{equation}
\det K^{AB} \;=\; \beta^4\left(1 - \tfrac{\alpha}{\beta}\right)
\;=\; \beta^3(\beta - \alpha) \;=\; -\beta^3(\alpha - \beta)\,,
\label{eq:detK_app}
\end{equation}
independent of $n^A(x)$.

From the densitized-metric identification $K^{AB} \propto
\sqrt{-g_{\mathrm{eff}}}\,g_{\mathrm{eff}}^{AB}$ and
$\det(\sqrt{-g}\,g^{\mu\nu}) = -g$ in four dimensions,
\begin{equation}
\det K^{AB} \;=\; -\,g_{\mathrm{eff}}\,,
\qquad
\sqrt{-g_{\mathrm{eff}}} \;=\;
\sqrt{\beta^3(\alpha-\beta)} \;=\; \frac{\beta^2}{c_*}\,,
\label{eq:sqrt_g_value_app}
\end{equation}
using $c_*^2 = \beta/(\alpha-\beta)$.

\paragraph{Vanishing variation on the constraint surface.} An
admissible variation $\delta n^A$ preserves the normalization
$n^A\delta_{AB}n^B=1$, hence $\delta_{AB} n^A \delta n^B = 0$. Differentiating
\eqref{eq:K_det_general_app},
\begin{equation}
\frac{\partial(\det K)}{\partial n^A}
\;=\; -2\alpha\beta^3\, \delta_{AB} n^B\,,
\end{equation}
hence
\begin{equation}
\delta(\det K)
\;=\; -2\alpha\beta^3\,(\delta_{AB} n^B \delta n^A) \;=\; 0\,,
\end{equation}
so $\delta\sqrt{-g_{\mathrm{eff}}} = 0$ on admissible variations, as
stated in Theorem~1(iii) of the main text.


\subsection{Jacobian cancellation identity for the classical
baseline}
\label{app:jacobian_cc}

The fundamental ECT action is $S_E[q] = \int_M d^4 x_E\,
\mathcal{L}_E(q,\partial q)$ with $q^a$ collectively denoting
$\Phi$, $n^A$, and auxiliary fields. Crucially, the emergent metric
$g_{\mathrm{eff}}^{AB}[q]$ is a functional of the condensate fields
rather than an independent dynamical variable; there is no variation
step ``by $g_{\mu\nu}$''.

Under the shift $\mathcal{L}_E \to \mathcal{L}_E + C$ for a
field-independent constant $C\in\mathbb{R}$:
(a) $\partial C/\partial q^a = \partial C/\partial(\partial_A q^a)
= 0$, so every Euler--Lagrange equation is unchanged; (b) the
classical solutions $q^a(x)$ are unchanged; (c) the emergent metric
$g_{\mathrm{eff}}^{AB}[q(x)]$ is unchanged on each classical
solution. Specialising to $C = V(\phi_0)$ establishes Theorem~1(iv).

The result is non-trivial in the ECT context precisely because
$g_{\mathrm{eff}}$ is emergent. In standard general relativity with
an independent $g_{\mu\nu}$, the same shift generates $C\sqrt{-g}$,
whose variation produces $\sim C g_{\mu\nu}\sqrt{-g}$ -- the standard
cosmological-constant source. That mechanism is unavailable here.


\subsection{FRW reduction and equation of state}
\label{app:frw_cc}

On an FRW background, spatial isotropy selects a single
long-wavelength ordered-branch degree of freedom $q(t)$. The
minimal homogeneous effective Lagrangian is
\begin{equation}
L_{\mathrm{IR}}[q; a] \;=\; a^3 \left[
\tfrac{1}{2} Z_q\, \dot q^2 - V_{\mathrm{IR}}(q) \right],
\label{eq:frw_lagrangian_app}
\end{equation}
with infrared stiffness $Z_q$ and residual infrared potential
$V_{\mathrm{IR}}(q)$. Crucially, $V_{\mathrm{IR}}(q)$ is the
dynamical dressing of the potential around the ordered-branch
minimum; it is not $V(\phi_0)$, which is decoupled by Theorem~1(iv).
This distinction resolves the apparent tension between the
decoupling theorem and the dark-energy picture used in
Sec.~\ref{sec:mp_late_cosmology}.

From \eqref{eq:frw_lagrangian_app},
\begin{align}
\rho_q &\;=\; \tfrac{1}{2} Z_q \dot q^2 + V_{\mathrm{IR}}(q)\,,
\\
p_q &\;=\; \tfrac{1}{2} Z_q \dot q^2 - V_{\mathrm{IR}}(q)\,,
\\
w_q &\;=\; \frac{r - 1}{r + 1}\,,
\qquad r \equiv
\frac{\tfrac12 Z_q \dot q^2}{V_{\mathrm{IR}}(q)}\,.
\end{align}
For $r \ll 1$ (slow-branch regime), $w_q \simeq -1 + 2r$,
reproducing the main-text formula $w_0 = -1 +
2\rho_{\mathrm{kin}}/(3\rho_{\mathrm{cond}})$ upon identification
$\rho_{\mathrm{kin}}=\tfrac12 Z_q\dot q^2$ and
$\rho_{\mathrm{cond}}=V_{\mathrm{IR}}(q)$.

\subsection{Dimensional audit of the infrared scaling}
\label{app:ir_audit_cc}

The observed $\rho_\Lambda$ has dimensions mass$^4$. After blocking
the direct UV channel (Theorem~1), the scale $M_{\mathrm{Pl}}^4$ is
not available as a direct source. Available infrared scales are the
emergent gravitational stiffness $M_{\mathrm{Pl}}^2$ (from the
Sakharov/Adler induced-gravity mechanism~\cite{Sakharov1967,Adler1982})
and the Hubble scale $H_0$ (only dimensionful IR curvature/time
scale).

Dimensionally consistent combinations of mass$^4$ from
$\{M_{\mathrm{Pl}}^2, H_0\}$ without directly invoking
$M_{\mathrm{Pl}}^4$ are $M_{\mathrm{Pl}}^2 H_0^2$, $M_{\mathrm{Pl}}
H_0^3$, and $H_0^4$. With $M_{\mathrm{Pl}}/H_0 \sim 10^{60}$, these
evaluate to $10^{-47}$, $10^{-107}$, $10^{-167}\,\mathrm{GeV}^4$
respectively; only the first matches observation at the
order-of-magnitude level.

\paragraph{Pseudo-Goldstone realization.} For a pseudo-Goldstone
$\theta$ of mass $m_G$ and decay constant $f_G$, the natural infrared
energy density is $\rho_G \sim m_G^2 f_G^2$, obtained from
$V_G(\theta) \sim m_G^2\theta^2$ at $\theta \sim f_G$. Using the ECT
values $m_G \sim H_0$ and $f_G \sim \phi_0 \sim M_{\mathrm{Pl}}$
(Sec.~\ref{subsec:goldstone_cc}),
\begin{equation}
\rho_G \;\sim\; m_G^2 f_G^2 \;\sim\; H_0^2 M_{\mathrm{Pl}}^2\,,
\end{equation}
reproducing the IR scaling with $c_\Lambda = O(1)$ via independent
inputs: $m_G$ comes from the pseudo-Goldstone (dark-matter-candidate)
sector, $f_G$ from the gravitational SSB scale. Neither is tuned
here.

\paragraph{Scope.} The estimate $\rho_G \sim H_0^2 M_{\mathrm{Pl}}^2$
is a consistency check, not a derivation. Upgrading it to a
prediction requires a first-principles derivation of $m_G$ from ECT
parameters (Open Problem~O4) and an analysis of the cosmological
evolution of the Goldstone condensate. The holographic bound of
Cohen--Kaplan--Nelson~\cite{CohenKaplanNelson1999} states
$\rho \lesssim M_{\mathrm{Pl}}^2/L_{\mathrm{IR}}^2$; the ECT estimate
saturates this bound with $L_{\mathrm{IR}}\sim H_0^{-1}$, which is an
independent phenomenological consistency check.

% =============================================================================
% END OF APPENDIX App_CC_technical
% =============================================================================